\chapter{Project Management and Development Evidence}
\label{chap:pm}

\section{Work Completed and Remaining}
Completed work is listed in Table~\ref{tab:i2-features} (Iteration~2) and
Section~\ref{sec:iter1-outcome} (Iteration~1). Remaining work is the FYP-2 plan
of Chapter~\ref{chap:fyp2}.

\section{Contribution of Each Team Member}
Table~\ref{tab:contributions} uses Git author identities on all branches
(non-merge commits). Commit counts are a repository metric, not a measure of
effort or quality: one commit may be one line or a whole module, and on
23~November 122 of Ali Sher's commits each added a single app file.

\begin{table}[H]\centering\footnotesize
\caption{Contribution evidence from Git history.}\label{tab:contributions}
\begin{tabularx}{\textwidth}{@{}L{2.6cm} L{3.2cm} l Y@{}}\toprule
\rowcolor{hdrfill}\textbf{Member} & \textbf{Git identity} & \textbf{Commits} & \textbf{Main areas evidenced}\\\midrule
Shahmeer Amir & \texttt{shah-paradox} & 47 & Live capture contracts, keyframes, gaps and buffer; the M1/M2 pass (gravity, relations, state runs, store, rooms, compaction, segment runner); Part~II plan; report and diagrams\\
Faizan ur Rehman Khan & \texttt{pez-Ox}, \texttt{Faizan Ur Rehman Khan} & 57 + 1 & Batch pipeline and R4DSG baseline; CI/CD and API; in-hand memory and Rerun; room hierarchy and registry; evaluation suite; demo harness; EgoLife study; tiered query protocol; SAM~3 benchmark\\
Ali Sher & \texttt{Ali Sher}, \texttt{Ali-Sher12} & 218 & Flutter redesign and search history; NLQ engine and \path{/ask}; Modal-served API; end-to-end runner; Patient and Caregiver Apps; enrolment; DINOv3 re-ID\\
\bottomrule\end{tabularx}\end{table}
\tbd{Each member to confirm their row and lead modules.}

\section{Milestones and Commit Evidence}
Figure~\ref{fig:commit-timeline} shows commits per week and member;
Figure~\ref{fig:branch-graph} shows the branches and the two tags. Work was
concentrated in two bursts: the mid-evaluation (weeks 40--42) and Iteration~2
(weeks 45--48).

\aegisfig{ch9-commit-timeline}{Commits per ISO week and author (non-merge).}{fig:commit-timeline}
\aegisfig{ch9-branch-graph}{Branches, merges and tags.}{fig:branch-graph}

\section{Repository and Version Control}
The repository uses a \texttt{src} layout (\path{src/aegis/}), YAML
configuration in \path{configs/}, tests in \path{tests/}, the apps in
\path{apps/}, and conventional commit messages. Work happens on feature
branches merged into \texttt{develop}; \texttt{main} receives tagged releases
(\texttt{v0.1.0-mid1}, \texttt{v0.2.0-mid2}). CI runs Ruff, pytest with the
coverage floor and Flutter checks. Appendix~\ref{app:repository} gives the
layout.

\section{Tools and Tool-Assisted Development}
\label{sec:tools}
VGGT-$\Omega$, SAM~2, SAM~3, DINOv3, YOLOE, FastAPI, Flutter, Modal, OpenCV and
Rerun are external components. AI assistants were used for coding help, for
drafting parts of this report in LaTeX and for diagram markup (PlantUML and
matplotlib scripts). The team reviewed all generated material and is responsible
for the architecture, algorithms, requirements, tests and claims.
\tbd{Per-member disclosure of which assistant was used for what, to be completed
by each member.}

\section{External Components and Dependencies}
\begin{table}[H]\centering\footnotesize
\caption{External dependencies.}\label{tab:external-dependencies}
\begin{tabularx}{\textwidth}{@{}L{4.6cm} Y l@{}}\toprule
\rowcolor{hdrfill}\textbf{Dependency} & \textbf{Use} & \textbf{Status}\\\midrule
Python 3.12, uv, NumPy, OpenCV, PyYAML & Runtime and video processing & Used\\
FastAPI, Uvicorn, httpx & API and tests & Used\\
SQLite (standard library) & Home store & Used\\
Modal, CUDA images, L4 GPU & Remote workers and API hosting & Used\\
VGGT-$\Omega$ & Poses and 3D points & Used\\
SAM~3, YOLOE, SAM~2 & Detection and masks & Used\\
DINOv3 & Re-identification & Used\\
Flutter/Dart, \texttt{speech\_to\_text} & Apps and on-device speech & Used\\
Rerun & Offline 3D viewer & Optional\\
Hugging Face Hub & Model weights (gated VGGT-$\Omega$, mirror fallback) & Used\\
Identity provider, encryption library & Authentication and encryption at rest & Planned\\
\bottomrule\end{tabularx}\end{table}

\section{Risks and Mitigation}
Likelihood (L) and impact (I) are scored 1--3; the score is L$\times$I.
\begin{table}[H]\centering\footnotesize
\caption{Scored risk register.}\label{tab:risk-register}
\begin{tabularx}{\textwidth}{@{}l Y l l l Y@{}}\toprule
\rowcolor{hdrfill}\textbf{ID} & \textbf{Risk} & \textbf{L} & \textbf{I} & \textbf{Score} & \textbf{Mitigation / status}\\\midrule
TC1 & Identity errors under clutter and occlusion & \tbd{2} & \tbd{3} & \tbd{6} & DINOv3 re-ID, debouncing; matching fix 61dda40\\
TC2 & Compaction loses useful detail & \tbd{1} & \tbd{3} & \tbd{3} & Verification before delete; last-known protected\\
TC3 & Motion blur and lighting & \tbd{2} & \tbd{2} & \tbd{4} & Keyframe gate\\
TC4 & GPU cost and latency & \tbd{2} & \tbd{2} & \tbd{4} & L4 \$0.069 per 5-min segment; CPU fallback\\
TC5 & Hallucinated answers & \tbd{1} & \tbd{3} & \tbd{3} & Closed-vocabulary rule engine; abstention codes\\
TC6 & High interaction load for older users & \tbd{2} & \tbd{2} & \tbd{4} & Object tiles, large targets; usability study planned\\
TC7 & Privacy and consent & \tbd{3} & \tbd{3} & \tbd{9} & Open: auth, encryption (Iteration~4)\\
TC8 & Limited validation & \tbd{3} & \tbd{2} & \tbd{6} & Measured/unmeasured separation; \texttt{demo05}\\
TC9 & Evaluation dataset & \tbd{2} & \tbd{2} & \tbd{4} & Ego4D NO\_GO; EgoLife used\\
TC10 & RGB-only spatial uncertainty & \tbd{2} & \tbd{2} & \tbd{4} & Gravity frame; no metric claims\\
TC11 & Room errors and cross-session identity & \tbd{2} & \tbd{2} & \tbd{4} & Registry, matching; metrics planned\\
TC12 & Room linking unstable across scale & \tbd{3} & \tbd{2} & \tbd{6} & Scale-relative linking (Iteration~3)\\
\bottomrule\end{tabularx}\end{table}

\section{Security and Privacy Considerations}
\label{sec:security}
The current prototype is not safe for real patient data. The API has no
authentication or authorisation; CORS allows any origin, method and header; the
writes \path{/rooms/rename}, \path{/enroll}, \path{/dataset/select} and
\path{/capture/pause} are open; and \texttt{GET /video/\{dataset\_id\}} serves
raw video to anyone with the URL. The home store, \texttt{home\_map.json} and
evidence crops are stored unencrypted. HTTPS is provided only by the Modal
endpoint. Speech is transcribed on the phone, so no audio leaves it. All testing
used team-recorded or public data. Iteration~4 addresses these gaps (P1.7, P1.8).

\section{Current Limitations}
Spatial and temporal accuracy, latency and usability are unmeasured; temporal
answers are not built; confidence is not calibrated; room linking is unstable
across scale; storage growth is projected from short clips; and no clinical
claim is made.
